% R588: EGA I, 0, 3.2.2; source-faithful draft, editorial review pending.
% Traducción: OpenAI Codex, GPT-6.1 Sol, esfuerzo Ultra.
\begin{env}[3.2.2]
\label{0.3.2.2-fr}
Para que el prehaz $\mathcal{F}'$ así definido sea un \emph{haz},
es necesario y suficiente que el prehaz $\mathcal{F}$ sobre $\mathfrak{B}$
satisfaga la condición~:

\smallskip
\noindent
\textup{(F\textsubscript{0})} \emph{Para todo recubrimiento $(U_\alpha)$ de
$U\in\mathfrak{B}$ por conjuntos $U_\alpha\in\mathfrak{B}$ contenidos en
$U$, y para todo objeto $T\in\boldsymbol{K}$, la aplicación que a cada
$f\in\operatorname{Hom}(T,\mathcal{F}(U))$ hace corresponder la familia
\[
  (\rho_{U_\alpha}^{U}\circ f)
  \in\prod_\alpha\operatorname{Hom}(T,\mathcal{F}(U_\alpha))
\]
es una biyección de $\operatorname{Hom}(T,\mathcal{F}(U))$ sobre el conjunto de las
familias $(f_\alpha)$ tales que
$\rho_V^{U_\alpha}\circ f_\alpha=\rho_V^{U_\beta}\circ f_\beta$ para todo
par de índices $(\alpha,\beta)$ y todo $V\in\mathfrak{B}$ tal que
$V\subset U_\alpha\cap U_\beta$}%
\footnote{Esto significa también que el par formado por $\mathcal{F}(U)$ y los
$\rho_\alpha=\rho_{U_\alpha}^{U}$ es \emph{solución del problema universal}
definido en (3.1.1) por los datos
$A_\alpha=\mathcal{F}(U_\alpha)$,
$A_{\alpha\beta}=\prod\mathcal{F}(V)$ (para los $V\in\mathfrak{B}$ tales que
$V\subset U_\alpha\cap U_\beta$) y
$\rho_{\alpha\beta}=(\rho_V''):\mathcal{F}(U_\alpha)\to
\prod\mathcal{F}(V)$ definido por la condición de que, para
$V\in\mathfrak{B}$, $V'\in\mathfrak{B}$, $W\in\mathfrak{B}$,
$V\cup V'\subset U_\alpha\cap U_\beta$, $W\subset V\cap V'$,
$\rho_W^V\circ\rho_V''=\rho_W^{V'}\circ\rho_{V'}''$.}.

La condición es evidentemente necesaria. Para demostrar que es suficiente,
consideremos primero una segunda base $\mathfrak{B}'$ de la topología de $X$,
\emph{contenida en $\mathfrak{B}$}, y demostremos que, si $\mathcal{F}''$
denota el prehaz deducido de la subfamilia
$(\mathcal{F}(V))_{V\in\mathfrak{B}'}$, $\mathcal{F}''$ es
\emph{canónicamente isomorfo} a $\mathcal{F}'$. En efecto, en primer lugar el
límite proyectivo (para $V\in\mathfrak{B}'$, $V\subset U$) de los morfismos
canónicos $\mathcal{F}'(U)\to\mathcal{F}(V)$ es un morfismo
$\mathcal{F}'(U)\to\mathcal{F}''(U)$ para todo abierto $U$. Si
$U\in\mathfrak{B}$, este morfismo es un isomorfismo, pues por hipótesis los
morfismos canónicos $\mathcal{F}''(U)\to\mathcal{F}(V)$ para
$V\in\mathfrak{B}'$, $V\subset U$, se factorizan en
\[
  \mathcal{F}''(U)\to\mathcal{F}(U)\to\mathcal{F}(V),
\]
y es inmediato comprobar que las composiciones de los morfismos
$\mathcal{F}(U)\to\mathcal{F}''(U)$ y
$\mathcal{F}''(U)\to\mathcal{F}(U)$ así definidos son las identidades. Así,
para todo abierto $U$, los morfismos
$\mathcal{F}''(U)\to\mathcal{F}''(W)=\mathcal{F}(W)$ para
$W\in\mathfrak{B}$ y $W\subset U$ satisfacen las condiciones que caracterizan el
límite proyectivo de los $\mathcal{F}(W)$ ($W\in\mathfrak{B}$, $W\subset U$), lo
que demuestra nuestra afirmación, habida cuenta de la unicidad de un límite proyectivo
salvo isomorfismo.

Dicho esto, sean $U$ un abierto cualquiera de $X$, $(U_\alpha)$ un recubrimiento
de $U$ por abiertos contenidos en $U$, y sea $\mathfrak{B}'$ la
subfamilia de $\mathfrak{B}$ constituida por los conjuntos

\oldpage[0\textsubscript{I}]{27}

de $\mathfrak{B}$ contenidos en algún $U_\alpha$; es claro que
$\mathfrak{B}'$ es todavía una base de la topología de $U$, por lo que
$\mathcal{F}'(U)$ (resp. $\mathcal{F}'(U_\alpha)$) es el límite proyectivo de los
$\mathcal{F}(V)$ para $V\in\mathfrak{B}'$ y $V\subset U$ (resp.
$V\subset U_\alpha$); el axioma (F) se comprueba entonces inmediatamente por la
definición del límite proyectivo.

Cuando se satisface (F\textsubscript{0}), diremos por abuso de lenguaje que
el prehaz $\mathcal{F}$ sobre la base $\mathfrak{B}$ es \emph{un haz}.
\end{env}
